\section{Cross-Domain Discussion: What Transfers and What Does Not}
\label{sec:discussion}

The empirical record now spans two structurally different markets:
LOB mid-price prediction \cite{solovev2026lob} and DeFi USDC lending-rate
forecasting (\Cref{sec:defi-experiment}).  We separate the discussion
into three claims at distinct layers of the cross-domain transfer
hypothesis --- architectural, predictive, and strategic --- and then
position our negative H1 verdict against the impossibility result of
\cite{halpern2024fair}.

\subsection{Architectural transfer succeeds}
\label{sec:disc-arch}

The DA-BiGRU-CNN of \Cref{sec:methodology-abstract} trains stably and
extracts non-trivial signal in both domains under \emph{identical}
network wiring, composite loss, and optimization schedule.  In LOB it
attains test weighted-Pearson $\LobGRUwPearson$, outperforming a
205-feature LightGBM by $\LobGRUOverLGBM$
\cite{solovev2026lob,ke2017lightgbm}; in DeFi it reaches validation
weighted-Pearson $0.747$ on Aave with directional accuracy $0.739$ and
a positive Compound $R^{2}=0.138$ (\Cref{sec:defi-forecaster}).  The
recipe that ports without modification is concrete: (i) identify a
domain-natural decomposition pair $(\mathcal{X}_A, \mathcal{X}_B)$
such that one subspace carries the deterministic backbone and the
other carries the residual stochastic component --- (price, volume) in
LOB, $(\varepsilon = r - f_{\text{kink}}(u), u+\text{ctx})$ in DeFi;
(ii) re-use the same dual-BiGRU stack with multi-scale Conv1d fusion;
(iii) re-use the same composite loss with weighted-Pearson as the
dominant term.  The \emph{architectural} transfer hypothesis ---
``the same template generalises by re-identifying the decomposition
pair'' --- is therefore supported by the validation diagnostics.  This
is the central methodological contribution and is independent of the
H1 outcome.

\subsection{Predictive transfer is partial}
\label{sec:disc-pred}

The transfer breaks one layer up.  Validation weighted-Pearson $0.747$
on Aave collapses to test $-0.07$, and directional accuracy falls from
$0.739$ on validation to $0.351$ on test --- a sub-chance value whose
sign is the empirical signature of a forecaster that has locked onto a
correlation pattern from the wrong regime.  The proximate cause is
documented: between 2025-Q3 and 2025-Q4 the Aave-pays-more share rises
from $39.9\%$ to $56.3\%$ and the median cross-protocol spread inverts
from $-0.22$\,pp to $+0.18$\,pp (\Cref{sec:bg-defi-mechanics}).  The
validation fold lives in the pre-inversion regime; the test fold
straddles the inversion and the 2026-Q1 reversion.  In-distribution
predictability is therefore not equivalent to robust short-horizon
predictability: the rate process exhibits regime structure on a
quarterly time-scale that the four-month test window happens to slice
adversarially.  The LOB benchmark \cite{solovev2026lob} did not exhibit
a comparable test-fold collapse, plausibly because LOB
microstructural dynamics are dominated by stationary
adverse-selection mechanics on millisecond horizons, whereas DeFi
lending rates inherit on-the-order-of-quarter regime persistence from
slow demand-side composition shifts.  This is not a falsification of
the architecture; it is a falsification of the implicit assumption
that high in-sample weighted-Pearson predicts out-of-sample
allocator edge.

\subsection{Strategic transfer fails on this window}
\label{sec:disc-strat}

Predictive MCDM under-performs reactive MCDM-EMA by $8$\,bp in
annualised net APY ($\PredictiveAPY$ vs $\EMAAPY$); the paired
monthly-Sharpe bootstrap (\Cref{sec:defi-h1}) returns
$\SharpeDelta$ with 95\% CI $[\BootstrapCILow,\,\BootstrapCIHigh]$
and $p=\BootstrapPvalue$, firmly inside the H0 acceptance region of
the pre-registered $\Delta\textsc{Sharpe}\ge 0.2$ test.  The
hysteresis-sweep ablation of \Cref{sec:defi-ablation} rules out the
most natural rebuttal --- ``the threshold is mis-tuned'' --- by
showing that the design $\theta=0.05$ sits at the Sharpe knee, with
fewer rebalances (and zero gas) actually \emph{raising} net APY in
this window.  The mechanism is therefore: a forecaster that lost
sign on test cannot drive a switching allocator into the right pool
ahead of the EMA-smoothed reactive comparator; the predictive
overhead consumes gas without obtaining timing alpha.  This is
exactly the negative-ensemble pattern that
\cite{solovev2026lob}~\S5.5 reported for LOB --- combining a strong
GRU with a 205-feature LightGBM \emph{degraded} test
weighted-Pearson by $0.05$--$1.6\%$ rather than improving it ---
generalised one layer up the stack: combining a forecaster with a
reactive allocator can underperform the allocator alone if the
forecaster transports out-of-distribution badly.  Two negative
results, in two structurally different domains, with the same
mechanistic explanation: \emph{the apparent complementarity of two
components is a property of the joint distribution they were
calibrated on, not a robust architectural fact.}

\subsection{Relation to the Halpern--Pass--Saraf impossibility}
\label{sec:disc-impossibility}

\cite{halpern2024fair} prove that no static interest-rate function can
correctly price a perpetual lending pool with early-repayment
optionality; any rate either over- or under-compensates lenders
relative to the option-implied premium.  The result constrains
\emph{contract design at the pricing layer}.  Our finding is at the
\emph{execution layer}, and it is empirical rather than theoretical:
short-horizon local predictability of the realised rate stream is
present in-distribution (validation weighted-Pearson $0.747$ on Aave)
but does not survive the 2025-Q3 $\to$ Q4 regime shift documented in
the data panel.  These two negative findings do not refute each other
--- one is a no-arbitrage statement about equilibrium fairness, the
other is an out-of-sample statement about the stationarity of the
short-horizon predictive structure --- and they do not jointly imply
that DeFi-lending forecasting is hopeless.  They jointly imply that
the empirical-predictability conjecture must be re-stated in a
regime-conditional form, with explicit handling of distribution shift
on the order of one fiscal quarter rather than a global predictability
claim across an 18-month panel.  The right response is methodological,
not pessimistic: longer test folds, online recalibration on regime
boundaries, and reporting test-window-conditional metrics
(\Cref{sec:limitations}, Finding~\#9).  Our H0-publishable verdict is
in this sense complementary to \cite{halpern2024fair}: it provides the
empirical data point that the impossibility-of-fair-pricing result
predicts will manifest as regime-persistent disequilibrium in the
realised rate stream.
